\documentclass[a4paper, twosided, 11pt]{scrartcl}

%% Language configuration
\usepackage{polyglossia}
\setdefaultlanguage[variant=swiss]{german}

%$ Specify where floatings go and make nicer looking tables
\usepackage{float}
\usepackage{array}
\usepackage{booktabs}
\usepackage{tabularx}

%% Set up font
\usepackage[T1]{fontenc}
\usepackage[usefilenames, DefaultFeatures={Ligatures=Common}]{plex-otf}
\renewcommand*\familydefault{\sfdefault}

%% Header and Footers
\usepackage[automark]{scrlayer-scrpage}
\ihead{Pflichtenheft}
\chead{}
\ohead{\headmark}

\ifoot{}
\cfoot{}
\ofoot{\pagemark}

%% Tikz to draw stuff
\usepackage{tikz}
\usepackage{tikzsymbols}
\usetikzlibrary{automata, positioning, arrows, calc, shapes.geometric, fit}

%% Access metadata
\usepackage{titling}

% Metadata
\title{Pflichtenheft}
\author{Naoki Pross, Joanne Robertz, Steven Wegmann, J\'er\^ome Roy}

% Document
\begin{document}
\begin{titlepage}
  \includegraphics[height=3cm]{pic/ost-logo}
  \begin{flushright}
    % TODO: OST logo
    \vspace{5cm}
    {\Huge \bfseries \thetitle} \\
    \vspace{5mm}
    {\LARGE Project: \textit{GlowingBroccoli}} \\
    \vspace{5mm}
    {\LARGE \bfseries Version 3.4}
  \end{flushright}
\end{titlepage}


\clearpage
\section*{\"Anderungsgeschichte}
\begin{tabularx}{\textwidth}{lrlX}
  \toprule
  \bfseries Datum & \bfseries Version & \bfseries Autor & \bfseries Beschreibung \\
  \midrule
  2020-10-05 & 1.0 & Robertz J.       & Dokument erstellt \\
  2020-10-19 & 2.0 & Robertz J.       & Einleitung und allgemeine Beschreibung bearbeiten \\
  2020-10-19 & 2.1 & Wegmann S.       & Funktionale Anforderungen bearbeiten \\
  2020-10-19 & 2.2 & Roy J.           & Dokument bearbeiten \\
  2020-10-19 & 2.3 & Pross N.         & Dokument bearbeiten \\
  2020-10-22 & 3.0 & Pross N.         & Dokument in \textrm{\LaTeX} gesetzt \\
  2020-10-23 & 3.1 & Pross N.         & Use Cases bearbeiten \\
  2020-10-24 & 3.2 & Roy J.           & Dokument bearbeiten \\
  2020-10-24 & 3.3 & Roy J., Pross N. & Use Cases bearbeiten \\
  2020-10-25 & 3.4 & Roy J., Pross N. & Typos und Fehlern korrigiert \\
  \bottomrule
\end{tabularx}

\vfill
{
  \renewcommand{\arraystretch}{2}
  \begin{tabularx}{\textwidth}{lp{.2\textwidth}X}
                   & \bfseries Datum & \bfseries Unterschrift \\
    Pross Naoki    & \hrulefill & \hrulefill \\
    Robertz Joanne & \hrulefill & \hrulefill \\
    Roy J\'er\^ome & \hrulefill & \hrulefill \\
    Wegmann Steven & \hrulefill & \hrulefill \\
  \end{tabularx}
}

\clearpage
\tableofcontents
\listoffigures
\listoftables

\clearpage

\section{Einleitung}
\subsection{Zweck}

Im vorliegenden Dokument sind die Anforderungen definiert, welche im Projekt
``GlowingBroccoli'' umgesetzt werden müssen. Es beschreibt den Auftrag zwischen
Auftraggeber und Auftragnehmer. Der Ausdruck Pflichtenheft ist hier im Sinne
der IEEE Recommended Practice for Software Requirements Specification.
ANSI/IEEE Std 830-1998 verwendet. Die dort definierte Requirements
Specification beinhaltet sowohl die Benutzeranforderungen (Lastenheft gemäss
DIN 69901-5) als auch Realisierungsvorgaben an die Entwicklungsgruppe
(Pflichtenheft gemäss DIN 69901-5).

\subsection{Produk\"uberblick}
\begin{figure}[H]
  \centering
  \begin{tikzpicture}
  \draw[lightgray, step = 5mm, dotted] (0,0) grid (10, 5);

  %% Food
  \node at (8, 2) (food) {};
  \draw[thick, fill = red!50!black] (food) rectangle +(-.5, -.5);

  %% Snake
  \node at (3, 4) (snake-head) {};
  \draw[thick, fill = yellow!70!black] (snake-head) rectangle +(-.5, -.5);
  \foreach \i in {1, ..., 3} {
    \draw[thick, fill = green!70!blue]
      ($ (snake-head) + (-.5 * \i, 0) $)
        node (snake-body-\i) {}
        rectangle +(-.5, -.5);
  }

  \foreach \i in {0, ..., 4} {
    \draw[thick, fill = green!70!blue] ($ (snake-body-3) + (0, -.5 * \i) $) rectangle +(-.5, -.5);
  }

  %% Arrows
  \draw[<-, very thick] (food) to[controls = +(45:1) and +(-90:1)] +(1, 3.5)
    node[above] {\textbf{Food}};

  \draw[<-, very thick] (snake-body-2.north) to[controls = +(90:1) and +(-90:1)] +(2, 1.5)
    node[above] {\textbf{Snake}};

  \draw[->, ultra thick, green!50!black] ($ (snake-head.east) - (0, .25) $) to +(1, 0)
    node[right] {};
\end{tikzpicture}

  \caption[Beispiel f\"ur grafische Oberfl\"ache in Use Case Play]{
    Beispiel dafür, wie die grafische Oberfläche f\"ur das Use-Case Play aussehen sollte.
  }
\end{figure}

Unser Team hat sich dafür entschieden das klassische ``Snake'' Spiel zu
programmieren.  Im Spiel sind Sie eine Schlange, deren Ziel es ist, sich auf
einer Karte zu bewegen und die zufällig auftauchende Nahrung zu essen. Wenn
eine Nahrungseinheit gegessen wird, dehnt sich die Schlange aus und die
Punkzahl wird erhöht.  Ziel des Spiels ist es, die höchstmögliche Punktzahl zu
erreichen.  Der Benutzer verliert sein Spiel, wenn der Kopf der Schlange einen
festen Gegenstand (die Wände der Spielkarte) berührt oder wenn der Kopf einen
Teil seines Körpers frisst.  In dieser Version der \emph{Snake} wird es
verschiedene Ebenen geben. Jede Stufe hat eine Siegpunktzahl zu erreichen, um
zur nächsten überzugehen. Wenn der Spieler die zuletzt vorgeschlagene Stufe
erreicht, hat er das Spiel gewonnen.


% Im Grunde geht es darum, durch bewegende Punkte, eine Schlange
% aufzubauen, die immer grösser wird, umso mehr sie isst.
% Als Abbruchbedingungen mit gefolgtem ``Game Over'' zählt:
% \begin{itemize}
%   \item Die Schlange berührt die Wand.
%   \item Die Schlange frisst sich selber, zum Beispiel weil sie so lang geworden
%     ist, dass es unmöglich ist sie im Spielfeld weiterhin agil zu bewegen.
% \end{itemize}

\subsection{Definitionen, Akronyme und Abk\"urzungen}

\begin{description}
  \item[Widget] Diese Umgebung wird innerhalb des Fensters aufgebaut und dient
    um einen dynamischen Verlauf, wie z.B. ein Spiel, aufzubauen.
  \item[Scoreboard] Liste der bisherigen Stücke nach Punktzahl geordnet.
\end{description}

\section{Allgemeine Beschreibung}
\subsection{System\"ubersicht}
\begin{figure}[h]
  \centering
  \begin{tikzpicture}[
    usecase/.style = {
      ellipse, draw, very thick, fill=gray!20,
      inner sep = 4pt, outer sep = 3pt
    },
    user/.pic = {
      \begin{scope}[line width = 1.5mm]
        \node[draw, circle, inner sep = .2cm] (head) {};
        %
        \draw (head.south) to +(0,-.2) node[inner sep = 0, outer sep = 0] (neck) {};
        \draw (neck.south) to +(0,-.5) node[inner sep = 0, outer sep = 0] (waist) {};
        %
        \draw (neck.center) to +(.5, -.2);
        \draw (neck.center) to +(-.5, -.2);
        %
        \draw (waist.center) to +(-.3, -.6);
        \draw (waist.center) to +(.3, -.6);
      \end{scope}
    },
  ]
  \begin{scope}
    \node[usecase] (menu) {Menu};
    \node[usecase, below = of menu] (game) {Play};
    \node[usecase, below = of game] (pause) {Pause};
    \node[usecase, right = of menu] (scoreboard) {Scoreboard};
    \node[usecase, right = of pause] (gameover) {Game Over};

    \begin{scope}[very thick]
      \draw[->, gray] (menu) to (game);
      \draw[<->, gray] (menu) to (scoreboard);
      \draw[<->, gray] (game) to (pause);
      \draw[<-, dashed] (game) to[out = 10, in = -120]
        node[pos = .7, above left] {\small extends} (scoreboard);
      \draw[<-, dashed] (game) to[out = -10, in = 120]
        node[pos = .7, below left] {\small extends} (gameover);
      \draw[->, gray] (gameover) to[out = 70, in = -70] (scoreboard);
    \end{scope}
  \end{scope}

  \node[very thick, draw, dashed, gray, fit = {
    ($(pause.south west) + (-1, -1)$)
    ($(scoreboard.north east) + (1, 1)$)
  }] (boundary) {};

  \node[above, gray] at (boundary.north) {\textbf{Snake Game}};

  \node[left = of boundary.west, outer sep = 5mm, xshift = -1cm] (player) {};
  \pic[at = {(player.center)}] {user};
  \node[below = of player] {Player};

  \draw[->, very thick] (player) to[controls = +(20:2) and +(180:2) ] (menu);
  \draw[->, very thick] (player) to[controls = +(0:2) and +(180:2) ] (game);
  \draw[->, very thick] (player) to[controls = +(-20:2) and +(180:2) ] (pause);

\end{tikzpicture}

  \caption{Use-case Diagramm}
\end{figure}

Das Spiel Snake verfügt über eine grafische Schnittstelle, mit der Sie
interagieren können.  Der Benutzer entscheidet, wann das Spiel durch Drücken
der Starttaste gestartet wird.  Zu Beginn des Spiels bewegt sich die Schlange
in eine bestimmte Richtung und es erscheint eine zufällige Portion Nahrung.  An
diesem Punkt wartet das Spiel auf Eingaben des Benutzers, um die Richtung zu
ändern. Der Benutzer kann die Richtung der Schlange mit den Pfeilen oder mit
den \texttt{WASD}-Tasten wählen. 

\subsection{Produktfunktionen} \label{sec:product-functions}
\begin{itemize}
  \item Das Spiel wird mittels einem Startknopf eröffnet.

  \item Die Schlange muss mittels den vier Pfeiltasten (oder \texttt{WASD})
    bewegt werden können.

  \item Das herumliegende Essen muss sich der Schlange anschliessen als Punkt,
    damit sie sich nach jedem Bissen auch vergrössert.

  \item ``Game Over'' muss erscheinen, wenn die Schlange aus einer bestimmten
    Abbruchsbedingung stirbt.

  \item Ein Scoreboard zeigt dem Benutzer die Ingame-Punkte der vorherigen Spiele an.
\end{itemize}

\subsection{Benutzereingenschaften}
Dieses Projekt soll Benutzergruppen ansprechen die gerne ihre Zeit in einen
lustigen Klassiker investieren möchten. Eine Altersbeschränkung gibt es hier
nicht. Auch Vorkenntnisse im Bereich der Informatik werden zum spielen nicht
gebraucht.

\subsection{Einschr\"ankungen}
Keite.

\subsection{Annahmen und Abh\"angigkeiten}
Die Software soll auf allen der drei Hauptplattformen: Windows 10, MacOS und
Linux laufen. Dies kann durch \texttt{C++} als Programmiersprache und das
plattformübergreifende Qt Framework erreicht werden. Von der User Hardware wird
erwartet, dass sie die Fähigkeit hat einfache Grafiken darzustellen und über
eine moderne grafische Umgebung verfügt.

\subsection{Priorisierung der Anforderungen}
Die Funktionen, welche in \S\ref{sec:product-functions} ersichtlich sind,
beschreiben die Muss-Anforderungen welche benötigt werden, um den minimalen
Zweck der Software zu erfüllen. Als Soll-Anforderungen wird die Implementation
von optisch schöneren Grafiken beschrieben.  Die Wunsch-Anforderung beinhaltet
das Einfügen von weiteren Funktionen im Spiel, wie zum Beispiel verschiedene
Früchte, mehrere \emph{Levels} oder die Implementation eines
``Scoreboards''.

\section{Externe Schnittstellen}
Es gibt keine speziellen externen Schnittstellen für diese Software. Sie
benötigen lediglich eine Tastatur und einen Bildschirm. 

\section{Funktionale Anforderungen}

\subsection{Actors}
\begin{table}[H]
  \centering
  \begin{tabular}{ll}
    \toprule
    Actors & Beschreibung \\
    \midrule
    Player & Die Eingaben des Users beeinflussen den Verlauf des Spiels. \\
    \bottomrule
  \end{tabular}
  \caption{Liste der Actors}
\end{table}

\subsection{Kurzbeschreibung der Use Cases}
\begin{table}[H]
  \centering
  \begin{tabular}{lp{.75\linewidth}}
    \toprule
    Use Case & Beschreibung \\
    \midrule
    % \texttt{Exit} & Beim betätigen der Systemtaste \texttt{E} verlässt man das Spiel
    % sofort und gelangt auf die Hauptseite, wo sich wiederum der Startknopf
    % befindet. \\None, except keyboard and display.

    Menu & Der Spieler kann einen Namen eingeben und wählen, ob er spielen oder
    die \emph{Scoreboard} sehen möchte. \\

    Play & Das Spiel \emph{Snake} spielen. \\

    Pause & Beim betätigen der Systemtaste \texttt{P} oder \texttt{ESC} wird
    das Spiel pausiert. Beim erneuten drücken, geht das Spiel weiter. \\

    Game Over & Der Spieler hat das Spiel verloren, sein Punktestand wird auf
    dem Bildschirm angezeigt. \\

    Scoreboard & Es wird eine Liste der früheren Spiele desselben oder eines
    anderen Spielers angezeigt. \\

    \bottomrule
  \end{tabular}
  \caption{Kurzbeschreibung der Use Cases}
\end{table}

\subsection{Use Case Menu}
\paragraph{Vorbedingungen} Es gibt eine grafische Umgebung und das Programm
muss laufen.
\paragraph{Nachbedingungen} Keine.
\paragraph{Nicht-Funktionale Anforderungen} Keine.
\paragraph{Haputszenario} Der Benutzer wird angefordert eine Entscheidung zu
treffen, das Spiel starten, ein Username w\"ahlen, das \emph{Scoreboard} sehen
oder das Programm schliessen.
\paragraph{Unterszenarien} Keine.
\paragraph{Fehlerszenarien} Es gibt keine grafische Umgebung.
\paragraph{Regeln} Keine.
\paragraph{Anmerkungen} Keine.
\paragraph{Beispiele} Keine.

\subsection{Use Case Play}
\paragraph{Vorbedingungen} Die Tastatur muss vorhanden sein.
\paragraph{Nachbedingungen} Keine.
\paragraph{Nicht-Funktionale Anforderungen} Keine.
\paragraph{Haputszenario} Das \emph{Snake} Spiel lauft.
\paragraph{Unterszenarien} Keine.
\paragraph{Fehlerszenarien} Es ist keine Username oder Tastatur vorhanden.
\paragraph{Regeln}
Die Schlange wird mit der Tastatur gesteuert. Die Schlange darf sich nicht aus
dem sichtbaren Bereich entfernen. Es ist der Schlange auch nicht erlaubt, ihren
eigenen Körper zu durchqueren. Wenn einer dieser zwei F\"alle geschieht, kommt es zu einem Game
Over.

\paragraph{Anmerkungen} Keine.
\paragraph{Beispiele} Keine.

\subsection{Use Case Pause}
\paragraph{Vorbedingungen} Es muss ein laufendes Spiel geben.
\paragraph{Nachbedingungen} Das Spiel ist pausiert.

\paragraph{Nicht-Funktionale Anforderungen} Bei Betätigung des Knopfes soll der
User mit möglichst wenig Verzögerung eine Pause der Spielpartie erzwingen
können, die er aber später weiterspielen kann.

\paragraph{Haputszenario} Ein Menu gibt die m\"oglichkeit zu fortfahren oder
zum Hauptmenu zu zur\"uckkehren.

\paragraph{Unterszenarien} Keine.

\paragraph{Fehlerszenarien} Es gibt kein laufendes Spiel.
\paragraph{Regeln}
Wenn das Spiel lauft, wird es unterbrochen. Wenn das Spiel bereits angehalten
ist, sollte es wieder aufgenommen werden. Es sollte auch möglich sein, das
Spiel zu beenden und zum Hauptmenü zurückzukehren.
\paragraph{Anmerkungen} Keine.
\paragraph{Beispiele} Keine.

\subsection{Use Case Game Over}
\paragraph{Vorbedingungen} Es muss ein laufendes Spiel geben.
\paragraph{Nachbedingungen} Das Spiel endet.
\paragraph{Nicht-Funktionale Anforderungen} Keine.
\paragraph{Haputszenario} Der Spieler wird informiert, dass er verloren hat.
Sein \emph{Score} (Punktzahl) wird auf den Bildschirm gezeigt.
\paragraph{Unterszenarien} Keine.
\paragraph{Fehlerszenarien} Keine.
\paragraph{Regeln} Der Spieler kann ein Neustart anfordern, das
\emph{Scoreboard} sehen oder zum Hauptmenü zurückzukehren.
\paragraph{Anmerkungen} Keine.
\paragraph{Beispiele} Die Schlange frisst sich selbst.

\subsection{Use Case Scoreboard}
\paragraph{Vorbedingungen} Programm muss laufen.
\paragraph{Nachbedingungen} Keine.
\paragraph{Nicht-Funktionale Anforderungen} Keine.
\paragraph{Haputszenario} Eine Liste der vorherige Spiele wird auf den
Bildschirm gezeigt.
\paragraph{Unterszenarien} Keine.
\paragraph{Fehlerszenarien} Keine.
\paragraph{Regeln} Keine.
\paragraph{Anmerkungen} Keine.
\paragraph{Beispiele} Keine.

\section{Sonstige Anforderungen}

\subsection{Entwurfseinschränkungen (Design Constraints)} Die grafische Qualität
des Spiels wird einfach und begrenzt sein, da wir die Grundfunktionen von Qt
zum Zeichnen der verschiedenen Komponenten verwenden.

\subsection{Wartbarkeit} Diese Software wurde für die aktuelle Version der
Betriebssysteme Linux, Windows 10 und Mac OS Catalina entwickelt. Es wird keine
Updates für neuere Betriebssysteme geben.

\subsection{Logging und Tracing} Der Benutzer muss sich nicht registrieren, um
das Spiel spielen zu können. Es gibt jedoch die Idee, vergangene Spiele lokal
zu speichern, so dass Sie die in verschiedenen Spielen erzielten Ergebnisse
vergleichen können.

\addcontentsline{toc}{section}{Referenzen}
\renewcommand{\refname}{Referenzen}
\begin{thebibliography}{2}
  \bibitem{ieee}
    IEEE Recommended Practice for Software Requirements Specification. ANSI/IEEE Std 830-1998
\end{thebibliography}

\end{document}
